\subsection{KRULL 整环中的高度 1 素理想}

定义 4. 设 A 是整环。若 A 的素理想 p 在 A 的非零素理想中是极小的，就称 p 的高度为 1。

我们也称理想 (0) 的高度为 0；因此，按定义，高度 \(\leqslant 1\) 的素理想等于 (0) 或者高度为 1。

下面我们将以一般方式定义素理想的高度。

定理 3. 设 A 是 Krull 整环，p 是 A 的整理想。p 是对应于某个极端除子的除性理想的充要条件是 p 是高度为 1 的素理想。

若 p 是对应于某个极端除子的除性理想，我们知道（第 4 节，命题 6 的推论 1）p 是素理想且 \(A_{p}\) 是离散赋值环；由于 \(A_{p}\) 除了 (0) 和 \(pA_{p}\) 外没有其他素理想，(0) 和 p 就是包含于 p 的 A 的唯一素理想（第 II 章，§ 3，第 1 节，命题 3）；因此 p 的高度为 1。反之，我们先证明每个 \(p \neq (0)\) 的 A 的素理想都包含一个对应于极端除子的除性素理想 q：因为 \(A_{p} \neq K\)，所以 \(A_{p}\) 是非空族（\(A_{t}\)）的本性赋值环的交（第 4 节，命题 6）；每个 \(A_{t}\) 都具有 \(A_{q_{t}}\) 的形式（第 4 节，命题 6 的推论 1），由 \(A_{p} \subset A_{q_{t}}\) 可推出 \(q_{t} \subset p\)。因此，若 p 的高度为 1，则 p = q，这说明 p 是对应于某个极端除子的除性理想。

推论 1. 在 Krull 整环中，每个非零素理想 m 都包含一个高度为 1 的素理想。若 m 的高度不是 1，则 div m = 0 且 A: m = A。

第一个断言已经在定理 3 的证明过程中得到。若 \(m\) 的高度不是 1，且 \(p\) 是一个高度为 1 的素理想并包含于 \(m\)，则 \(p \subset \tilde{m}\) 且 \(p \neq \tilde{m}\)；由于 \(\text{div } p\) 是极端的，必有 \(\text{div } m = \text{div } \tilde{m} = 0\)；从而 \(\text{div}(A: m) = 0\)，又因为 \(A: m\) 是除性的（第 1 节，命题 1），所以 \(A: m = A\)。

推论 2. 设 A 是 Krull 整环，K 是其分式域，v 是 K 上在 A 上为正的赋值，p 是 \(x \in A\) 中满足 \(v(x) > 0\) 的元素的集合。若素理想 p 的高度为 1，则 v 等价于 A 的一个本性赋值。

设 \(B\) 是 \(v\) 的赋值环，\(m\) 是其理想。则 \(m \cap A = p\)，从而 \(A_p \subset B\)。现在 \(A_p\) 是离散赋值环（定理 3 以及命题 6 的推论 1）。由于 \(p \neq (0)\)，\(B \neq K\)，因而 \(B = A_p\)（第 VI 章，§ 4，第 5 节，命题 6）。

定理 4. 设 A 是整环，M 是其高度为 1 的素理想集合。A 为 Krull 整环的充要条件是满足以下性质：

\begin{enumerate}
\def\labelenumi{(\roman{enumi})}
\item
  对所有 \(\mathfrak{p} \in \mathbf{M}\)，\(\mathbf{A}_{\mathfrak{p}}\) 都是离散赋值环。
\item
  A 是各个 \(\mathbf{A}_{\mathfrak{p}}\) 的交，其中 \(\mathfrak{p} \in \mathbf{M}\)。
\item
  对所有 \(x \neq 0\)（属于 \(\mathbf{A}\)），指标 \(\mathfrak{p} \in \mathbf{M}\) 中满足 \(x \in \mathfrak{p}\) 的至多有有限个。
\end{enumerate}

此外，与 \(A_{p}\) 对应的、其中 \(p \in M\) 的赋值正是 A 的本性赋值。

这些条件显然充分。其必要性立即由第 4 节的定理 3、命题 6 的推论 1 以及 A 的本性赋值满足第 3 节定义 3 的条件这一事实得出。

命题 10. 设 A 是整闭的诺特整环，a 是 A 的整理想。以下条件等价：

\begin{enumerate}
\def\labelenumi{(\alph{enumi})}
\item
  a 是除性的；
\item
  与 A/a 相伴的素理想都具有高度1。
\end{enumerate}

回忆一下，若 \(a = \bigcap_{i=1}^{n} q_i\) 是 a 的既约准素分解，且 \(p_i\) 表示对应于 \(q_i\) 的素理想，则与 A/a 相伴的素理想正是 \(p_i\)（第 IV 章，§2，第 3 节，命题 4）。于是 (a) 蕴含 (b) 由第 4 节的命题 8 得出。反之，若在上述记号下 \(p_i\) 的高度为 1，则 \(A_{p_i}\) 是离散赋值环（定理 4）；现在，\(q_i = q_i A_{p_i} \cap A\)（第 IV 章，§2，第 1 节，命题 3）；记 \(v_i\) 为对应于 \(p_i\) 的本性赋值，则存在整数 \(n_i\)，使得 \(q_i\) 是 \(x \in A\) 的集合，其中 \(v_i(x) \geqslant n_i\)；这说明 \(q_i\) 是除性的（第 4 节，命题 5），从而 a 也是除性的。

